\documentclass[10pt,a4paper]{article}
\usepackage[utf8]{inputenc}
\usepackage[english]{babel}
\usepackage{amsmath}
\usepackage{amsfonts}
\usepackage{amssymb}
\usepackage{xcolor}
\usepackage{graphicx}

\newcommand{\prob}[1]{\mathrm{Prob}\lbrace #1\rbrace}
\newcommand{\power}[1]{\mathrm{PSD} ( #1 ) }
\newcommand{\laplace}[1]{\mathcal{L}\lbrace #1\rbrace}
\newcommand{\expected}[1]{\mathrm{E}[ #1 ]}
\newcommand{\variance}[1]{\mathrm{V}[ #1 ]}
\newcommand{\heaviside}[1]{\mathrm{\theta}\left( #1 \right)}

\begin{document}

  \section{ General Renewal Process }

  \subsection{ Asymptotic Distribution }

  Consider the following signal
  \begin{equation}
    u(t) = \chi_k, \mbox{ with } t \in [T_k, T_{k+1}], \mbox{ and } \chi_{k+1} = -\chi_k
  \end{equation}
  The times $T_i$ are considered independent and identically distributed random
  variables with expected mean value $\expected{T_i} = \mu$, variance
  $\variance{T_i} = \sigma$, probability density $q(t)$ and cumulative
  probability density $Q(t)$, while $N(t)$ is the number of events inside the
  interval $[0,t)$. We also consider the waiting time to the n-th event.
  \begin{equation}
    G_n = T_1 + T_2 + \cdots + T_n
  \end{equation}
  such that
  \begin{equation}
    \prob{G_n > t} = \prob{N(t) < n}
  \end{equation}
  It can be shown that the distribution of $N(t)$ as $t\to\infty$ is
  asymptotically normal with expected mean value $\expected{N(t)} = t/\mu$ and
  variance $\variance{N(t)} = \sigma^2 t/\mu^3$.

  \subsection{ Relation between the autocorrelation function and the power spectrum}

  In applied mathematics, the Wiener–Khinchin theorem, states that the
  autocorrelation function of a wide-sense-stationary random process has a
  spectral decomposition given by the power spectrum of that process
  \begin{equation}
    \power{\omega} = \int_{-\infty}^{\infty} C(t) \exp{(-i\omega t)} ~dt, \mbox{ with }
    \omega = 2\pi f
  \end{equation}
  where the autocorrelation function reads as
  \begin{equation}
    C(t) \equiv \frac{1}{T_f}\int_0^{T_f} u(t')u(t+t') dt'
  \end{equation}
  The value $u(t')u(t+t')$ may be written as $\chi_k\chi_l$, where $t'\in [T_l, T_{l+1})$
  and $t+t'\in [T_k, T_{k+1})$. The idea is to approximate the autocorrelation
  function using the probability distribution $q(\tau)$ by considering only the
  time intervals where $\chi_k\chi_l > 0$,
  \begin{equation}
    C(t) \approx \frac{1}{\mu} \int_t^{T_f} (\tau - t) q(\tau) d\tau
  \end{equation}

  Let us suppose that $q(\tau) = K \tau^{-\beta}$ where $K$ is some
  proportionality constant. Integration of equation (6) leaves
  \begin{equation}
    C(t) \approx \frac{K}{\mu} \frac{(2-\beta)(1-\beta)}{(2-\beta)(1-\beta)} t^{2-\beta}, \mbox{ with }
    \beta > 2
  \end{equation}
  \begin{align*}
    \power{f} &\sim \int_{-\infty}^{\infty} t^{2-\beta} \exp{(-i\omega t)} ~dt
    \\
                  &\sim f^{\beta-3} \int_{-\infty}^{\infty} u^{2-\beta} \exp{(-iu)} ~du
  \end{align*}
  Suggesting that the power spectrum decays as $\alpha = \beta - 3$ for $2 < \beta$.
  A similar expression may be recovered for $1 < \beta < 2$.

  \paragraph{Question: } a) Does $q(\tau) \sim \tau^{-\beta}$ fits to any of our
  observations? b) How about $\power{f} \sim f^{-\alpha}$? If so, for which
  frequency range?


  \subsection{ A model for the probability distribution $q(t)$}

  Let us consider the survival function as the probability that an event will
  survive beyond a time $t$.
  \begin{equation}
    s(t) = \prob{T_i > t} = 1 - Q(t)
  \end{equation}
  Also consider the hazard function as the instantaneous failure rate
  \begin{equation}
    h(t) = \lim_{dt\to 0} \frac{\prob{t < T_i < t + dt}}{dt} \frac{1}{s(t)} = \frac{q(t)}{s(t)}
  \end{equation}
  also expressed as a cumulative hazard function
  \begin{equation}
    \Lambda(t) = \int_0^t h(t) dt = -\log{(s(t))}
  \end{equation}

  The example given by [Herault] consists in two attractors in the phase-space:
  one attractor A, corresponding to the value $u=+1$, and another attractor B,
  corresponding to the value $u=-1$. The system moves from one attractor to the
  other with a given probability $\Gamma(t)$, which corresponds precisely to
  the hazard function $h(t)$, while the probability function $P(t)=\Gamma(t)f(t)$
  corresponds precisely to the expression above $q(t)= s(t)h(t)$.

  \smallskip

  Let us assume a constant hazard rate $h(t)=\lambda$, meaning the probability
  of observing a future event does not event on the time passed from the previous
  event, i.e. the system has no memory of previous events. In such case,
  \begin{align}
    h(t) &= \lambda
    \\
    \Lambda(t) &= \int_0^t h(t) dt = \lambda t
    \\
    s(t) &= \exp{(-\Lambda(t))} = \exp{(-\lambda t)}
    \\
    q(t) &= s(t)h(t) = \lambda \exp{(-\lambda t)}
  \end{align}
  one obtains the Exponential distribution.

  Now, let us assume a hazard rate $h(t)$ that is uniformly distributed between
  two values $\lambda_1$ and $\lambda_2$,
  \begin{align}
    h(t) &=
    \begin{cases}
      \frac{1}{\lambda_2 - \lambda_1},& \mbox{ for } \lambda_1 < t < \lambda_2
      \\
      0 ,&\mbox{ otherwise }
    \end{cases}
    \\
    \Lambda(t) &= \int_0^t h(t) dt =
    \frac{(t - \lambda_1) - (t - \lambda_2)\heaviside{t - \lambda_2}}{\lambda_2 - \lambda_1} \heaviside{t - \lambda_1}
    \\
    s(t) &= \exp{(-\Lambda(t))} = \exp{\left(
    -\frac{(t - \lambda_1) - (t - \lambda_2)\heaviside{t - \lambda_2}}{\lambda_2 - \lambda_1} \heaviside{t - \lambda_1}
    \right)}
    \\
    q(t) &= s(t)h(t) =
    \begin{cases}
      \frac{1}{\lambda_2 - \lambda_1}
      \exp{\left(
      -\frac{t - \lambda_1}{\lambda_2 - \lambda_1}
      \right)},& \mbox{ for } \lambda_1 < t < \lambda_2
      \\
      0 ,&\mbox{ otherwise }
    \end{cases}
  \end{align}

  Now, [Herault] considers a system with some form of memory given by
  \begin{align}
    h(t) &= \frac{\nu}{\theta + t}
    \\
    \Lambda(t) &= \int_0^t h(t) dt = -\log{\left(\frac{\theta}{\theta + t}\right)^\nu}
    \\
    s(t) &= \exp{(-\Lambda(t))} = \left(\frac{\theta}{\theta + t}\right)^\nu
    \\
    q(t) &= s(t)h(t) = \frac{\nu \theta^\nu }{(\theta + t)^{\nu+1}}
  \end{align}

  In his thesis, [Herault] considers the case where moving from A to B, and
  moving from B to A, have different probabilities
  \begin{equation}
    h_{A\to B}(t) = \frac{\nu_1}{\theta + t}, \mbox{ and }
    h_{B\to A}(t) = \frac{\nu_2}{\theta + t}
  \end{equation}
  and uses Monte Carlo simulations to study the case with $\nu_1=\nu_2$,
  $\nu_1<\nu_2$, and $\nu_1\ll\nu_2$. Two different behaviours are observed: the
  proposed $\alpha+\beta=3$ for symmetric configurations, and $\alpha-\beta=1$
  for the strongly asymmetric configurations, the latter being more relevant for
  burst intermittent dynamics.

  %
  % \paragraph{What other possibilities?}
  % Let us consider the hazard function for the Weibull distribution
  % \begin{align}
  %   h(t) &= \gamma x^{\gamma-1}
  %   \\
  %   \Lambda(t) &= \int_0^t h(t) dt = t^\gamma
  %   \\
  %   s(t) &= \exp{(-\Lambda(t))} =  \exp{(-t^\gamma)}
  %   \\
  %   q(t) &= s(t)h(t) = \gamma x^{\gamma-1} \exp{(-t^\gamma)}
  % \end{align}
  % If $\gamma<1$, the probability of observing future events decreases with time
  % and {\it vice versa}. Note that $\gamma=1$ corresponds to a Poisson process.
  %
  %  Let us consider the following hazard rate
  %  \begin{align}
  %    h(t) &= \frac{\frac{1}{t\sigma}\phi\left(\frac{\ln t}{\sigma}\right)}{\Phi\left(-\frac{\ln t}{\sigma}\right)}
  %    \\
  %    \Lambda(t) &= \int_0^t h(t) dt = -\ln{\left(1 - \Phi\left(\frac{\ln t}{\sigma}\right) \right)}
  %    \\
  %    s(t) &= \exp{(-\Lambda(t))} =  1 - \Phi\left(\frac{\ln t}{\sigma}\right)
  %    \\
  %    q(t) &= s(t)h(t) = \frac{1}{t\sigma}\phi\left(\frac{\ln t}{\sigma}\right)
  %  \end{align}
  %  where $\phi(t)=\frac{e^{-\frac{(x-\mu)^2}{2\sigma^2}}}{\sigma\sqrt{2\pi}}$ is the probability density function of the normal
  %  distribution and $\Phi(t)$ is the corresponding cumulative density function.
  %


\end{document}
